Our results suggest that, for this model and this set of confusable
classes, prediction errors are not generally explained by the model
attending to the wrong part of the image. In every misclassified example we
inspected, SHAP attribution fell on the animal itself rather than the
background. Instead, errors appear linked to whether the specific texture
cue that separates two similar classes -- coat pattern, facial marking --
was visible and well-lit in the source image. When that cue is present (as
in the correctly classified leopard image), the model gets the label right;
when it is obscured by lighting, crop, or shadow (as in the husky and
jaguar images), the model falls back to a visually similar but incorrect
class.

\subsection{Limitations}
\label{sec:limitations}

\textbf{Explainer resolution.} We use \texttt{shap.Explainer} with an
\texttt{Image} masker and \texttt{max\_evals=100}, which produces a coarse
superpixel grid (roughly $4\times4$ regions per image). This keeps runtime
practical (about 30 seconds per image on CPU) but means fine-grained
features -- ear size relative to head, individual whisker patterns -- are
not resolved as distinct regions. The kit fox case in
Section~\ref{sec:correct} illustrates this: we can say attribution falls on
the head, but not whether it specifically falls on the ears. A finer masker
grid or a higher \texttt{max\_evals} budget would sharpen this at the cost
of substantially longer runtime.

\textbf{Blur masker as a perturbation choice.} SHAP's Image masker
approximates "removing" a region by blurring it, not by deleting it. Blur
is a reasonable stand-in for occlusion, but it does not fully remove
texture or color information the way a black/grey patch or inpainting
would, so attribution magnitudes should be read as relative importance
within our specific perturbation scheme rather than an absolute measure of
feature necessity.

\textbf{Single image per class.} Each class in our test set is represented
by exactly one image. This is enough to demonstrate the method and to
surface concrete, inspectable examples of correct and incorrect
predictions, but it does not let us distinguish a class-level tendency from
an artifact of that one photo's lighting, pose, or crop. The husky and
jaguar errors we highlight could plausibly disappear on a better-lit photo
of the same class.

\textbf{Small, curated class set.} We selected 10 groups of classes because
they are visually similar by inspection, not through a systematic
similarity metric over the model's full confusion matrix. Our 23\% error
rate on this set is therefore a statement about deliberately hard cases, not
an estimate of the model's overall error rate.

\textbf{Top-1 attribution only.} We compute SHAP attribution for the
model's single top-predicted label. We do not inspect attribution for the
"correct" label when the model is wrong, which could additionally show
whether the true class was a close runner-up or effectively invisible to
the model.
